%!TEX root = ../../main.tex
%% ===============================================================
%% 5.3 특수 킬 마커와 동점 처리
%% 파일: chapters/05_label/03_special_kills_ties.tex
%% 상위: chapters/05_label/chapter.tex
%% 정의 라벨: sec:label-special
%% 참조 라벨: sec:disc-audit
%% 포함 객체: -
%% 인용: -
%% 상태: 초안 (docs/OUTLINE.md 참고)
%% ===============================================================

\section{특수 킬 마커와 동점 처리}
\label{sec:label-special}

Riot API는 하나의 사망에 대해 \texttt{CHAMPION\_KILL}과, 조건에 따라 \texttt{CHAMPION\_SPECIAL\_KILL} (첫 킬, 다중 킬, 에이스)을 \emph{함께} 발행한다. 후자는 별도의 킬이 아니라 같은 킬에 대한 표식이므로, 본 논문은 특수 마커에 대해서는 보너스 $s(u)$만 더하고 킬 가치 $1.0$은 짝이 되는 \texttt{CHAMPION\_KILL}에만 부여한다.\footnote{초기 구현은 특수 마커를 두 번째 킬로 계수하였고 사전 감사에서 수정되었다(제 \ref{sec:disc-audit} 절).} $S(U) = 0$인 동점은 시드 고정 무작위 배정(\texttt{random})을 기본으로 하며, 제외(\texttt{drop})나 한쪽 배정 옵션을 둔다. 동점은 전체의 극히 일부이며 라벨 균형은 약 50/50이다.
